\usepackage[a4paper,top=25mm,bottom=25mm,left=27mm,right=27mm,
            headheight=17pt,headsep=8mm,footskip=12mm]{geometry}
\usepackage{xcolor}
\usepackage{graphicx}
\usepackage{tikz}
\usetikzlibrary{mindmap,arrows.meta,positioning}
\usepackage{listings}
\usepackage{tcolorbox}
\tcbuselibrary{breakable}
\usepackage{booktabs,tabularx,longtable}
\usepackage{multicol}
\usepackage{caption}
\usepackage{fancyhdr}
\usepackage{xurl}
\usepackage{needspace}
\usepackage[unicode,bookmarksopen=true]{hyperref}

\newcommand{\BookTitlePrefix}{Python 3}
\newcommand{\BookTitleText}{核心编程与实战笔记}
\newcommand{\BookTitle}{\BookTitlePrefix\space\BookTitleText}
\newcommand{\BookSubtitle}{基于 PyCharm 与现代工程化实践}
% 封面与 PDF 作者字段共用以下设置。
\newcommand{\BookAuthor}{雪落长安无声}
\newcommand{\BookVersion}{v1.0}
% 默认使用编译日期。出版时可在这里填写固定日期。
\newcommand{\BookDate}{\today}

\definecolor{PythonBlue}{HTML}{3776AB}
\colorlet{PyBlue}{PythonBlue}
\definecolor{PythonYellow}{HTML}{FFD43B}
\definecolor{Ink}{HTML}{172B40}
\definecolor{Muted}{HTML}{4C6278}
\definecolor{DeepBlue}{HTML}{254E70}
\definecolor{CoverDark}{HTML}{122236}
\definecolor{EditorDark}{HTML}{1A3048}
\definecolor{EditorLine}{HTML}{35536E}
\definecolor{CodeBg}{HTML}{F1F6FB}
\definecolor{CodeLine}{HTML}{AFC6DB}
\definecolor{CodeKeyword}{HTML}{245B8A}
\definecolor{CodeString}{HTML}{75388B}
\definecolor{CodeComment}{HTML}{35685E}
\definecolor{PitfallRed}{HTML}{B43B49}

\setmainfont{TeX Gyre Heros}
\setsansfont{TeX Gyre Heros}
\setmonofont{lmmono10-regular.otf}[
  BoldFont=lmmonolt10-bold.otf,
  ItalicFont=lmmono10-italic.otf,
  BoldItalicFont=lmmonolt10-boldoblique.otf,
  Scale=MatchLowercase
]
\defaultCJKfontfeatures{}
\setCJKmainfont{FandolHei-Regular.otf}[
  BoldFont=FandolHei-Bold.otf,ItalicFont=FandolKai-Regular.otf
]
\setCJKsansfont{FandolHei-Regular.otf}[
  BoldFont=FandolHei-Bold.otf,ItalicFont=FandolKai-Regular.otf
]
\setCJKmonofont{FandolFang-Regular.otf}[
  BoldFont=FandolHei-Bold.otf,ItalicFont=FandolKai-Regular.otf
]

\setlength{\parindent}{0pt}
\setlength{\parskip}{0.65em plus 0.15em minus 0.1em}
\linespread{1.16}
\raggedbottom
\setcounter{tocdepth}{1}
\setcounter{secnumdepth}{2}

\ctexset{
  part={
    name={第,篇},
    number=\chinese{part},
    format=\centering\Huge\bfseries\sffamily\color{PyBlue},
    hang=false,
    beforeskip=92pt,
    afterskip=64pt,
    aftername={\par\nobreak\vspace{14pt}},
    aftertitle={\par\nobreak\vspace{12pt}{\color{PythonYellow}\rule{32mm}{2pt}}},
    pagestyle=empty
  },
  chapter={
    name={第\enspace,\enspace 章},
    number=\arabic{chapter},
    format=\raggedright\sffamily\bfseries\Huge\color{PythonBlue},
    aftername=\hspace{0.8em},
    beforeskip=24pt,
    afterskip=18pt
  },
  section={
    format=\raggedright\sffamily\bfseries\Large\color{Ink},
    beforeskip=18pt,
    afterskip=8pt
  }
}

\pagestyle{fancy}
\fancyhf{}
\fancyhead[L]{\small\ttfamily\color{PythonBlue}>>> \BookTitlePrefix}
\fancyhead[R]{\small\color{Muted}\BookTitleText}
\fancyfoot[C]{\small\ttfamily\color{Muted}\thepage}
\renewcommand{\headrulewidth}{0.4pt}
\renewcommand{\headrule}{%
  {\color{CodeLine}\hrule width\headwidth height\headrulewidth}%
}
\fancypagestyle{plain}{%
  \fancyhf{}%
  \fancyfoot[C]{\small\ttfamily\color{Muted}\thepage}%
  \renewcommand{\headrulewidth}{0pt}%
}

\hypersetup{
  pdftitle={\BookTitle},
  pdfauthor={\BookAuthor},
  pdfsubject={\BookSubtitle},
  colorlinks=true,
  linkcolor=DeepBlue,
  urlcolor=CodeKeyword
}
% 目录沿用自动生成的页码与链接，章号只显示阿拉伯数字。
\ctexset{
  chapter/tocline={
    \CTEXifname{\protect\numberline{\thechapter}}{\protect\numberline{}}#2
  }
}

\makeatletter
% 将篇标题和它下面的第一章一起换栏。
\newbox\BookPartTocBox
\let\BookChapterTocEntry\l@chapter
\renewcommand*{\l@part}[2]{%
  \box\BookPartTocBox
  \global\setbox\BookPartTocBox=\vbox{%
    \begingroup
      \hypersetup{linkcolor=PyBlue}%
      \color{PyBlue}%
      \parindent\z@
      \rightskip\@pnumwidth
      \parfillskip-\@pnumwidth
      \large\bfseries
      \noindent #1\nobreak\hfill
      \hb@xt@\@pnumwidth{\hss#2}\par
    \endgroup
  }%
}
\renewcommand*{\l@chapter}[2]{%
  \ifvoid\BookPartTocBox
    \BookChapterTocEntry{#1}{#2}%
  \else
    \addpenalty{-\@highpenalty}%
    \addvspace{1.3em}%
    \vbox{\box\BookPartTocBox\BookChapterTocEntry{#1}{#2}}%
    \global\setbox\BookPartTocBox=\box\voidb@x
  \fi
}
\renewcommand{\tableofcontents}{%
  \begingroup
    \ctexset{chapter={%
      format=\centering\sffamily\mdseries\Huge\color{PythonBlue},
      beforeskip=0pt,
      afterskip=26pt
    }}%
    \chapter*{\contentsname}%
    \@mkboth{\contentsname}{\contentsname}%
    \hypersetup{linkcolor=Ink}%
    \setlength{\columnsep}{11mm}%
    \setlength{\multicolsep}{0pt}%
    \setlength{\parskip}{0pt}%
    \renewcommand{\@pnumwidth}{2.1em}%
    \renewcommand{\@tocrmarg}{2.65em}%
    \renewcommand*{\l@section}{\addvspace{0.25em}\@dottedtocline{1}{0.75em}{2.45em}}%
    \begin{multicols}{2}
      \@starttoc{toc}%
      \box\BookPartTocBox
      \global\setbox\BookPartTocBox=\box\voidb@x
    \end{multicols}%
  \endgroup
}
\makeatother

\urlstyle{same}
\newcommand{\code}[1]{\texttt{\detokenize{#1}}}

% 代码块默认使用 Python 样式。ctex 在 XeLaTeX 下为 listings 接入中文支持。
\lstdefinestyle{python}{
  language=Python,
  basicstyle=\small\ttfamily\color{Ink},
  keywordstyle=\bfseries\color{CodeKeyword},
  stringstyle=\color{CodeString},
  commentstyle=\itshape\color{CodeComment},
  deletekeywords={print,exec},
  morekeywords={async,await,nonlocal},
  emph={print,len,range,enumerate,zip,open,int,str,list,dict,set,tuple},
  emphstyle=\color{DeepBlue},
  numbers=left,
  numberstyle=\scriptsize\ttfamily\color{Muted},
  stepnumber=1,
  numbersep=10pt,
  frame=single,
  rulecolor=\color{CodeLine},
  backgroundcolor=\color{CodeBg},
  framesep=8pt,
  framexleftmargin=22pt,
  xleftmargin=30pt,
  xrightmargin=9pt,
  columns=fullflexible,
  keepspaces=true,
  showstringspaces=false,
  tabsize=4,
  breaklines=true,
  upquote=true,
  aboveskip=12pt,
  belowskip=12pt
}
\lstset{style=python}

% 提示框均可跨页。
\newcommand{\NoteSquare}[1]{\textcolor{#1}{\rule{0.75em}{0.75em}}}
\tcbset{
  python note/.style={
    breakable,
    coltitle=Ink,
    coltext=Ink,
    fonttitle=\sffamily\bfseries,
    boxrule=0.85pt,
    arc=2mm,
    left=10pt,
    right=10pt,
    top=7pt,
    bottom=7pt,
    before skip=12pt,
    after skip=12pt
  }
}
\newtcolorbox{myunderstanding}{
  python note,
  title={\NoteSquare{PythonBlue}\enspace 核心直觉},
  colframe=PythonBlue,
  colback=PythonBlue!8!white,
  colbacktitle=PythonBlue!17!white
}
\newtcolorbox{pycharmpractice}{
  python note,
  title={\NoteSquare{PythonYellow}\enspace PyCharm 实战},
  colframe=PythonYellow,
  colback=PythonYellow!7!white,
  colbacktitle=PythonYellow!25!white
}
\newtcolorbox{pitfall}{
  python note,
  title={\NoteSquare{PitfallRed}\enspace 踩坑记录},
  colframe=PitfallRed,
  colback=PitfallRed!5!white,
  colbacktitle=PitfallRed!12!white
}

% 提示框图标使用 TikZ 矢量路径。
\newcommand{\SnakeIcon}{%
  \tikz[baseline=-0.55ex,x=1em,y=1em,line cap=round]{%
    \draw[line width=1.15pt]
      (0.12,0.18) -- (0.60,0.18)
      arc[start angle=-90,end angle=90,radius=0.15]
      -- (0.30,0.48) arc[start angle=-90,end angle=-270,radius=0.14]
      -- (0.72,0.76);
    \fill (0.73,0.76) circle[radius=0.10];
    \fill[PythonYellow!28!white] (0.75,0.80) circle[radius=0.023];
  }%
}
\newcommand{\BulbIcon}{%
  \tikz[baseline=-0.55ex,x=1em,y=1em,line cap=round,line join=round]{%
    \draw[line width=0.75pt]
      (0.29,0.29) -- (0.29,0.41)
      .. controls (0.02,0.64) and (0.21,0.89) .. (0.43,0.89)
      .. controls (0.65,0.89) and (0.84,0.64) .. (0.57,0.41)
      -- (0.57,0.29) -- cycle;
    \draw[line width=0.75pt] (0.34,0.14) -- (0.52,0.14);
  }%
}
\newcommand{\BugIcon}{%
  \tikz[baseline=-0.55ex,x=1em,y=1em,line cap=round]{%
    \draw[line width=0.75pt] (0.44,0.40) ellipse[x radius=0.22,y radius=0.28];
    \draw[line width=0.75pt]
      (0.44,0.64) -- (0.44,0.13)
      (0.32,0.65) -- (0.23,0.83) (0.56,0.65) -- (0.65,0.83)
      (0.25,0.52) -- (0.06,0.60) (0.63,0.52) -- (0.82,0.60)
      (0.22,0.38) -- (0.04,0.38) (0.66,0.38) -- (0.84,0.38)
      (0.26,0.25) -- (0.10,0.17) (0.62,0.25) -- (0.78,0.17);
  }%
}

\newtcolorbox{pycharmtip}{
  python note,
  title={\SnakeIcon\enspace PyCharm 快捷键},
  colframe=PythonYellow,
  colback=PythonYellow!7!white,
  colbacktitle=PythonYellow!28!white
}
\newtcolorbox{pythonictip}{
  python note,
  title={\BulbIcon\enspace Pythonic 写法},
  colframe=PythonBlue,
  colback=PythonBlue!10!white,
  colbacktitle=PythonBlue!22!white
}
\newtcolorbox{pybug}{
  python note,
  title={\BugIcon\enspace 常见报错与避坑},
  colframe=PitfallRed,
  colback=PitfallRed!5!white,
  colbacktitle=PitfallRed!13!white
}

% 思维导图样式可用于总览和章首。每张图可单独指定分支角度与距离。
\tikzset{
  book mindmap/.style={
    root concept/.style={
      minimum size=30mm,text width=24mm,font=\sffamily\bfseries\large
    },
    level 1 concept/.style={
      minimum size=23mm,text width=17mm,font=\sffamily\bfseries\small,
      level distance=38mm,sibling angle=72
    },
    level 2 concept/.style={
      minimum size=19mm,text width=16mm,font=\sffamily\footnotesize,
      level distance=27mm,sibling angle=52
    },
    mindmap,
    every node/.style={concept,align=center,text=Ink,inner sep=1mm},
    concept color=PythonBlue
  },
  python flow step/.style={
    rectangle,rounded corners=2mm,draw=PythonBlue,fill=CodeBg,
    text=Ink,align=center,font=\sffamily\small,
    minimum height=9mm,inner sep=5pt
  },
  python flow arrow/.style={-{Stealth[length=2.4mm]},draw=PythonBlue,thick}
}

\DeclareCaptionFont{captionink}{\color{Muted}}
\captionsetup[figure]{font={small,captionink},labelfont=bf,
                      justification=centering,singlelinecheck=false,position=bottom,
                      skip=9pt}

% 图像使用完整项目中的矢量 PDF。缺少图文件时使用同主题 TikZ 简图。
\newcommand{\BookImage}[2]{%
  \IfFileExists{figures/#1.pdf}{%
    \includegraphics[width=#2]{figures/#1.pdf}%
  }{%
    \resizebox{#2}{!}{\csname BookFallback#1\endcsname}%
  }%
}
\newcommand{\CaseStep}[2]{\par\smallskip{\bfseries\color{Muted}#1}\enspace #2\par}
\ctexset{
  part={
    break={\clearpage},
    format=\raggedright\Huge\bfseries\sffamily\color{PyBlue},
    fixskip=true,beforeskip=0pt,afterskip=8pt,hang=true,
    aftername=\hspace{0.65em},
    aftertitle={\par\nobreak{\color{PythonYellow}\rule{\linewidth}{1.5pt}}},
    pagestyle=plain
  },
  chapter={
    break={\clearpage},
    format=\raggedright\sffamily\bfseries\LARGE\color{PyBlue},
    fixskip=true,beforeskip=0pt,afterskip=9pt,pagestyle=plain
  },
  section={beforeskip=13pt,afterskip=6pt}
}
\newcommand{\BookPartIntro}{\vfill}
% 在篇页内排导读，并用弹性间距将寄语推至版心底部。
\newcommand{\BookPart}[4]{%
  \begingroup
  \renewcommand{\BookPartIntro}{%
    \vspace{14mm}%
    {\large\bfseries\sffamily\color{PyBlue}本篇导读\par}%
    \smallskip
    {\large\color{Ink}#2\par}%
    \medskip
    #3\par
    \vfill
    {\centering\large\sffamily\color{PyBlue}#4\par}%
    \vspace{5mm}%
  }%
  \part{#1}%
  \endgroup
}
\makeatletter
\renewcommand{\@endpart}{\par\BookPartIntro\clearpage}
\renewcommand{\backmatter}{\clearpage\@mainmatterfalse}
\makeatother
\renewcommand{\topfraction}{0.92}
\renewcommand{\bottomfraction}{0.85}
\renewcommand{\textfraction}{0.06}
\renewcommand{\floatpagefraction}{0.85}
\setcounter{topnumber}{4}
\setcounter{bottomnumber}{3}
\setcounter{totalnumber}{6}
\setlength{\textfloatsep}{10pt plus 2pt minus 2pt}
\setlength{\intextsep}{9pt plus 2pt minus 2pt}
\setlength{\LTpre}{6pt}
\setlength{\LTpost}{6pt}

\expandafter\def\csname BookFallbackinterpreter\endcsname{%
\begin{tikzpicture}[node distance=8mm]
\node[python flow step,text width=29mm] (a) {项目设置};
\node[python flow step,text width=29mm,right=of a] (b) {运行配置};
\node[python flow step,text width=29mm,right=of b] (c) {终端路径};
\draw[python flow arrow] (a) -- (b);
\draw[python flow arrow] (b) -- (c);
\node[above=5mm of b,font=\sffamily\bfseries\large,text=PythonBlue] {解释器核对};
\end{tikzpicture}
}

\expandafter\def\csname BookFallbackreferences\endcsname{%
\begin{tikzpicture}[node distance=8mm]
\node[python flow step,text width=29mm] (a) {名字};
\node[python flow step,text width=29mm,right=of a] (b) {外层容器};
\node[python flow step,text width=29mm,right=of b] (c) {共享的内层对象};
\draw[python flow arrow] (a) -- (b);
\draw[python flow arrow] (b) -- (c);
\node[above=5mm of b,font=\sffamily\bfseries\large,text=PythonBlue] {对象引用};
\end{tikzpicture}
}

\expandafter\def\csname BookFallbacktype-info\endcsname{%
\begin{tikzpicture}[node distance=8mm]
\node[python flow step,text width=29mm] (a) {编辑器推断};
\node[python flow step,text width=29mm,right=of a] (b) {断点现场};
\node[python flow step,text width=29mm,right=of b] (c) {运行时类型};
\draw[python flow arrow] (a) -- (b);
\draw[python flow arrow] (b) -- (c);
\node[above=5mm of b,font=\sffamily\bfseries\large,text=PythonBlue] {两层观察};
\end{tikzpicture}
}

\expandafter\def\csname BookFallbacktraceback\endcsname{%
\begin{tikzpicture}[node distance=8mm]
\node[python flow step,text width=29mm] (a) {异常末行};
\node[python flow step,text width=29mm,right=of a] (b) {文件和行号};
\node[python flow step,text width=29mm,right=of b] (c) {最小复现};
\draw[python flow arrow] (a) -- (b);
\draw[python flow arrow] (b) -- (c);
\node[above=5mm of b,font=\sffamily\bfseries\large,text=PythonBlue] {读回溯};
\end{tikzpicture}
}

\expandafter\def\csname BookFallbacklist-preview\endcsname{%
\begin{tikzpicture}[node distance=8mm]
\node[python flow step,text width=29mm] (a) {原列表};
\node[python flow step,text width=29mm,right=of a] (b) {切片新列表};
\node[python flow step,text width=29mm,right=of b] (c) {独立追加};
\draw[python flow arrow] (a) -- (b);
\draw[python flow arrow] (b) -- (c);
\node[above=5mm of b,font=\sffamily\bfseries\large,text=PythonBlue] {列表预览};
\end{tikzpicture}
}

\expandafter\def\csname BookFallbackset-difference\endcsname{%
\begin{tikzpicture}[node distance=8mm]
\node[python flow step,text width=29mm] (a) {请求名单};
\node[python flow step,text width=29mm,right=of a] (b) {已有名单};
\node[python flow step,text width=29mm,right=of b] (c) {待同步差集};
\draw[python flow arrow] (a) -- (b);
\draw[python flow arrow] (b) -- (c);
\node[above=5mm of b,font=\sffamily\bfseries\large,text=PythonBlue] {集合比较};
\end{tikzpicture}
}

\expandafter\def\csname BookFallbackgrouping\endcsname{%
\begin{tikzpicture}[node distance=8mm]
\node[python flow step,text width=29mm] (a) {原始事件};
\node[python flow step,text width=29mm,right=of a] (b) {状态计数};
\node[python flow step,text width=29mm,right=of b] (c) {按状态分组};
\draw[python flow arrow] (a) -- (b);
\draw[python flow arrow] (b) -- (c);
\node[above=5mm of b,font=\sffamily\bfseries\large,text=PythonBlue] {同源日志};
\end{tikzpicture}
}

\expandafter\def\csname BookFallbackscope\endcsname{%
\begin{tikzpicture}[node distance=8mm]
\node[python flow step,text width=29mm] (a) {当前函数};
\node[python flow step,text width=29mm,right=of a] (b) {外层和模块};
\node[python flow step,text width=29mm,right=of b] (c) {内置名字};
\draw[python flow arrow] (a) -- (b);
\draw[python flow arrow] (b) -- (c);
\node[above=5mm of b,font=\sffamily\bfseries\large,text=PythonBlue] {名字查找};
\end{tikzpicture}
}

\expandafter\def\csname BookFallbackdecorator\endcsname{%
\begin{tikzpicture}[node distance=8mm]
\node[python flow step,text width=29mm] (a) {定义原函数};
\node[python flow step,text width=29mm,right=of a] (b) {完成包装};
\node[python flow step,text width=29mm,right=of b] (c) {调用时执行};
\draw[python flow arrow] (a) -- (b);
\draw[python flow arrow] (b) -- (c);
\node[above=5mm of b,font=\sffamily\bfseries\large,text=PythonBlue] {装饰器};
\end{tikzpicture}
}

\expandafter\def\csname BookFallbacktransaction\endcsname{%
\begin{tikzpicture}[node distance=8mm]
\node[python flow step,text width=29mm] (a) {开始与执行};
\node[python flow step,text width=29mm,right=of a] (b) {成功提交};
\node[python flow step,text width=29mm,right=of b] (c) {异常回滚};
\draw[python flow arrow] (a) -- (b);
\draw[python flow arrow] (b) -- (c);
\node[above=5mm of b,font=\sffamily\bfseries\large,text=PythonBlue] {事务};
\end{tikzpicture}
}

\expandafter\def\csname BookFallbackconcurrency\endcsname{%
\begin{tikzpicture}[node distance=8mm]
\node[python flow step,text width=29mm] (a) {启动任务};
\node[python flow step,text width=29mm,right=of a] (b) {等待与协作};
\node[python flow step,text width=29mm,right=of b] (c) {恢复处理};
\draw[python flow arrow] (a) -- (b);
\draw[python flow arrow] (b) -- (c);
\node[above=5mm of b,font=\sffamily\bfseries\large,text=PythonBlue] {任务交接};
\end{tikzpicture}
}

\expandafter\def\csname BookFallbackhttp\endcsname{%
\begin{tikzpicture}[node distance=8mm]
\node[python flow step,text width=29mm] (a) {方法与路径};
\node[python flow step,text width=29mm,right=of a] (b) {服务处理};
\node[python flow step,text width=29mm,right=of b] (c) {状态与内容};
\draw[python flow arrow] (a) -- (b);
\draw[python flow arrow] (b) -- (c);
\node[above=5mm of b,font=\sffamily\bfseries\large,text=PythonBlue] {HTTP 往返};
\end{tikzpicture}
}

\expandafter\def\csname BookFallbackarchitecture\endcsname{%
\begin{tikzpicture}[node distance=8mm]
\node[python flow step,text width=29mm] (a) {窗口或接口};
\node[python flow step,text width=29mm,right=of a] (b) {业务规则};
\node[python flow step,text width=29mm,right=of b] (c) {数据存储};
\draw[python flow arrow] (a) -- (b);
\draw[python flow arrow] (b) -- (c);
\node[above=5mm of b,font=\sffamily\bfseries\large,text=PythonBlue] {业务分层};
\end{tikzpicture}
}

\expandafter\def\csname BookFallbackarray\endcsname{%
\begin{tikzpicture}[node distance=8mm]
\node[python flow step,text width=29mm] (a) {原数组};
\node[python flow step,text width=29mm,right=of a] (b) {共享视图};
\node[python flow step,text width=29mm,right=of b] (c) {独立副本};
\draw[python flow arrow] (a) -- (b);
\draw[python flow arrow] (b) -- (c);
\node[above=5mm of b,font=\sffamily\bfseries\large,text=PythonBlue] {数组与切片};
\end{tikzpicture}
}

\expandafter\def\csname BookFallbackquality\endcsname{%
\begin{tikzpicture}[node distance=8mm]
\node[python flow step,text width=29mm] (a) {原始四行};
\node[python flow step,text width=29mm,right=of a] (b) {去重三行};
\node[python flow step,text width=29mm,right=of b] (c) {有效两项};
\draw[python flow arrow] (a) -- (b);
\draw[python flow arrow] (b) -- (c);
\node[above=5mm of b,font=\sffamily\bfseries\large,text=PythonBlue] {数据核对};
\end{tikzpicture}
}

\expandafter\def\csname BookFallbackpackage\endcsname{%
\begin{tikzpicture}[node distance=8mm]
\node[python flow step,text width=29mm] (a) {源码与配置};
\node[python flow step,text width=29mm,right=of a] (b) {构建 wheel};
\node[python flow step,text width=29mm,right=of b] (c) {新环境运行};
\draw[python flow arrow] (a) -- (b);
\draw[python flow arrow] (b) -- (c);
\node[above=5mm of b,font=\sffamily\bfseries\large,text=PythonBlue] {交付链路};
\end{tikzpicture}
}

\expandafter\def\csname BookFallbackdiagnosis\endcsname{%
\begin{tikzpicture}[node distance=8mm]
\node[python flow step,text width=29mm] (a) {代码与数据};
\node[python flow step,text width=29mm,right=of a] (b) {文件与环境};
\node[python flow step,text width=29mm,right=of b] (c) {外部服务};
\draw[python flow arrow] (a) -- (b);
\draw[python flow arrow] (b) -- (c);
\node[above=5mm of b,font=\sffamily\bfseries\large,text=PythonBlue] {排错路线};
\end{tikzpicture}
}

% 图与表的题注均放在内容下方，长题注也按居中方式排版。
\captionsetup*[table]{font={small,captionink},labelfont=bf,
  justification=centering,singlelinecheck=false,position=bottom,skip=9pt}
